\chapter{Fundamentação teórica}
\label{cap:fundamentacao-teorica}

Este capítulo apresenta os conceitos que fundamentam o desenvolvimento e a avaliação do sistema proposto. Inicialmente, discutem-se as secas no Semiárido e a gestão de reservatórios. Em seguida, são abordados a simulação hidrológica, o balanço hídrico, a evaporação, as regras de operação, os sistemas com múltiplos reservatórios, os sistemas de suporte à decisão, as ferramentas de simulação aplicadas a sistemas de reservatórios e a verificação de modelos computacionais.

\section{Secas e variabilidade hidrológica no Semiárido}
\label{sec:secas-semiarido}

A seca é um fenômeno natural caracterizado pela redução prolongada da disponibilidade de água em relação às condições esperadas para determinada região. Sua manifestação pode ser observada em diferentes componentes do ciclo hidrológico, como precipitação, umidade do solo, vazões e armazenamento. De acordo com \citeonline{deNys2016}, o fenômeno deve ser diferenciado da aridez, que representa uma condição climática permanente, e da escassez hídrica, que também pode resultar do desequilíbrio entre disponibilidade e uso da água.

No Semiárido nordestino, a precipitação concentra-se em poucos meses e apresenta forte variação entre anos consecutivos. Como consequência, as vazões dos rios intermitentes são muito irregulares: anos de afluência abundante podem ser seguidos por sequências plurianuais de afluências reduzidas. Essa irregularidade, combinada a taxas de evaporação elevadas, faz com que a capacidade de armazenamento e a forma de utilizá-lo sejam determinantes para o atendimento das demandas \cite{campos2005}.

Os impactos das secas não dependem exclusivamente da intensidade ou da duração do déficit hídrico. Eles também são influenciados pela exposição da população, pela dependência econômica de atividades sensíveis ao clima, pela infraestrutura disponível e pela capacidade institucional de preparação e resposta. \citeonline{marengo2022} destacam que a combinação entre irregularidade das chuvas, degradação do solo e vulnerabilidade socioeconômica amplia os efeitos das secas no Nordeste.

Historicamente, as respostas às secas no Nordeste passaram por ações emergenciais, implantação de infraestrutura hidráulica e, mais recentemente, instrumentos de gestão integrada. \citeonline{campos2015} observa que as políticas públicas contribuíram para formar uma infraestrutura capaz de fornecer água às populações, sem eliminar a vulnerabilidade. O evento de 2012 a 2018 apresentou duração e severidade excepcionais e evidenciou a permanência de práticas reativas \cite{pontesFilho2020}. A gestão proativa, que envolve monitoramento, avaliação de riscos e planejamento antecipado de medidas, requer instrumentos capazes de apoiar decisões em diferentes estágios da seca \cite{deNys2016}.

\section{Gestão e operação de reservatórios}
\label{sec:gestao-reservatorios}

Os reservatórios constituem uma das principais estruturas empregadas para regularizar a disponibilidade hídrica. Por meio do armazenamento, parte dos volumes afluentes em períodos favoráveis pode ser utilizada posteriormente, reduzindo a diferença temporal entre oferta e demanda. Para \citeonline{campos2005}, o dimensionamento e a operação dessas estruturas estão associados à regularização de vazões e à avaliação do desempenho do sistema diante da variabilidade das afluências.

A operação envolve conflitos entre objetivos. Retirar mais água no presente aumenta o atendimento imediato, mas reduz o estoque disponível para meses futuros e eleva o risco de colapso em uma seca prolongada. Manter o reservatório cheio reduz esse risco, mas aumenta as perdas por evaporação e por vertimento. Além disso, diferentes usos, como abastecimento humano, irrigação e indústria, competem pela mesma água e podem ter prioridades distintas \cite{labadie2004,loucksVanBeek2017}.

A gestão de recursos hídricos não se limita à determinação da capacidade física das obras. \citeonline{camposStudart2003} destacam que a gestão contemporânea incorpora gestão da demanda, planejamento de bacias e participação dos usuários. Para a operação de reservatórios, isso implica considerar não apenas quanto pode ser armazenado, mas como e quando a água deve ser disponibilizada.

\section{Simulação hidrológica}
\label{sec:simulacao-hidrologica}

Simular um sistema consiste em representar, por meio de um modelo, a evolução de seu estado ao longo do tempo sob regras e entradas especificadas. Na operação de reservatórios, o modelo de simulação recebe uma sequência de afluências e de demandas, aplica uma regra de operação e calcula, intervalo a intervalo, o armazenamento, as retiradas, as perdas e as falhas de atendimento \cite{loucksVanBeek2017}.

Uma prática usual consiste em submeter o sistema à série histórica de vazões. Embora a sequência futura não seja igual à observada, a série histórica contém secas e cheias efetivamente ocorridas e permite examinar como uma configuração teria respondido a esses eventos. Os resultados obtidos dessa forma respondem a perguntas do tipo ``o que aconteceria se'', e não constituem previsões da operação futura \cite{loucksVanBeek2017,mcmahonAdeloye2005}.

É necessário distinguir simulação, previsão e otimização. A previsão estima valores futuros de variáveis, como a afluência dos próximos meses. A otimização busca, entre alternativas admissíveis, aquela que melhor atende a uma função objetivo. A simulação, por sua vez, calcula as consequências de uma alternativa previamente especificada. A revisão de \citeonline{yeh1985} mostra que simulação e otimização são abordagens complementares: a simulação permite representar regras com detalhamento e avaliar alternativas propostas pelo analista, enquanto a otimização orienta a busca por soluções. Este trabalho concentra-se na simulação.

\section{Balanço hídrico de reservatórios}
\label{sec:balanco-hidrico}

O comportamento de um reservatório pode ser representado pelo princípio da conservação de massa. Em escala mensal, uma formulação geral do balanço hídrico é dada pela Equação \ref{eq:balanco-hidrico}.

\begin{equation}
    V_{t+1} = V_t + I_t + T_t - R_t - E_t - S_t
    \label{eq:balanco-hidrico}
\end{equation}

em que $V_t$ e $V_{t+1}$ representam os volumes armazenados no início e no final do mês $t$; $I_t$ é o volume afluente; $T_t$ é o saldo de transferências; $R_t$ é a retirada efetiva; $E_t$ é a perda por evaporação; e $S_t$ é o vertimento. O volume final deve permanecer entre os limites mínimo e máximo fisicamente admissíveis, de modo que o excesso em relação à capacidade é contabilizado como vertimento e a insuficiência de volume reduz a retirada efetiva \cite{mcmahonAdeloye2005}.

Quando o volume disponível é insuficiente para atender integralmente à retirada solicitada, ocorre uma falha de atendimento e um déficit. A frequência, a duração e a magnitude dessas falhas ao longo da série caracterizam o desempenho da operação. \citeonline{hashimoto1982} propõem descrever esse desempenho por critérios de confiabilidade, associada à frequência de estados satisfatórios, resiliência, associada à rapidez de recuperação, e vulnerabilidade, associada à severidade das falhas.

\section{Evaporação e curvas cota-área-volume}
\label{sec:evaporacao-cav}

Em regiões semiáridas, a evaporação a partir do espelho d'água pode representar parcela expressiva do balanço hídrico. O volume evaporado em um mês é o produto da lâmina evaporada pela área da superfície livre. Como essa área varia com o armazenamento, ela é obtida pela relação cota-área-volume do reservatório, determinada por levantamento topográfico ou batimétrico \cite{mcmahonAdeloye2005}.

Como o armazenamento se altera ao longo do mês, a área no início do intervalo não representa adequadamente toda a superfície exposta. Uma aproximação usual consiste em adotar a área média entre o início e o fim do mês. Como a área final depende do volume final, que por sua vez depende da evaporação, o cálculo é implícito; ele pode ser resolvido de forma aproximada por uma estimativa preliminar do volume final, seguida do cálculo da área correspondente e da recomputação da evaporação com a área média. Nesse procedimento, se a lâmina $e_t$ for expressa em metros e a área $A$ em quilômetros quadrados, o produto $e_t\,A$ resulta diretamente em hectômetros cúbicos. Essa forma de tratar a evaporação em modelos de balanço mensal é descrita por \citeonline{mcmahonAdeloye2005} e \citeonline{loucksVanBeek2017}.

\section{Regras de operação e curvas-guia}
\label{sec:regras-operacao}

Uma regra de operação estabelece como o sistema deve responder às condições observadas. Entre as variáveis que podem orientar essa resposta estão o armazenamento, o período do ano, a demanda e a expectativa de afluências. Em situações de seca, regras previamente definidas permitem substituir decisões improvisadas por medidas graduais e transparentes \cite{deNys2016}.

As curvas-guia, também denominadas curvas de níveis meta, dividem o armazenamento em faixas mensais associadas a estados operacionais. Neste trabalho são considerados os estados Normal, Alerta, Seca e Seca Severa. Acima do primeiro limite, o reservatório opera em condição Normal. À medida que o volume atravessa os limites de Alerta, Seca e Seca Severa, percentuais progressivos de redução da demanda, denominados racionamento, podem ser aplicados. Essa redução antecipada das retiradas, conhecida na literatura como \textit{hedging}, aceita déficits menores no presente para reduzir a probabilidade de déficits severos no futuro \cite{shihReVelle1994}. Como os limites variam ao longo do ano, a regra pode exigir mais volume no início da estação seca, quando a perspectiva de recarga é menor, e aceitar volumes menores ao final da estação chuvosa.

A permanência em um estado corresponde à proporção de meses nos quais o armazenamento é classificado naquela faixa. Esse indicador descreve a frequência de enquadramento operacional e não equivale à garantia de atendimento de uma demanda: uma política pode manter o reservatório em estado restritivo e, ainda assim, atender integralmente à demanda reduzida prevista para essa faixa. A efetividade da regra depende da posição das curvas e dos percentuais de racionamento: limites excessivamente baixos podem ser acionados tarde demais, enquanto limites excessivamente altos impõem restrições desnecessárias \cite{shihReVelle1994}.

\section{Operação de sistemas com múltiplos reservatórios}
\label{sec:multiplos-reservatorios}

Quando dois ou mais reservatórios são operados de forma integrada, a regra deve definir também a responsabilidade de cada unidade e as condições de transferência. A operação conjunta pode aproveitar a complementaridade entre reservatórios com capacidades, afluências e perdas diferentes, mas aumenta o número de decisões e de interações a representar \cite{labadie2004}.

Em uma configuração em série, um reservatório pode fornecer água a outro, por exemplo por meio de uma adutora, de um canal ou da liberação no leito do rio. A transferência costuma ser acionada quando o reservatório receptor atinge um volume de referência, denominado gatilho, e é limitada pela disponibilidade do fornecedor. Em uma configuração em paralelo, os reservatórios podem atender a uma demanda comum, e a regra define qual unidade é responsável pelo atendimento em cada momento; quando a unidade prioritária atinge seu gatilho, a responsabilidade pode ser deslocada para outra unidade \cite{lundGuzman1999}.

Essas representações são simplificações da infraestrutura real. Em intervalo mensal, por exemplo, uma transferência é tratada como instantânea, e perdas em trânsito, tempos de percurso e restrições hidráulicas precisam ser explicitamente representados ou admitidos como hipóteses \cite{loucksVanBeek2017}.

\section{Sistemas de suporte à decisão}
\label{sec:ssd}

Sistemas de suporte à decisão são ferramentas destinadas a auxiliar a análise de problemas nos quais dados, modelos e julgamento humano precisam ser utilizados conjuntamente. \citeonline{sprague1980} caracteriza esses sistemas pelo apoio a decisões semiestruturadas ou não estruturadas, pela interação com o usuário e pela possibilidade de combinar informações e modelos analíticos. Sua finalidade não é automatizar a decisão, mas ampliar a capacidade de investigar alternativas.

Na estrutura proposta por \citeonline{sprague1980}, um SSD reúne uma base de dados, uma base de modelos e mecanismos de comunicação com o usuário. A base de dados armazena as informações necessárias à análise; a base de modelos representa o comportamento do problema; e a interface permite configurar entradas e interpretar resultados. No campo da gestão de águas, \citeonline{souzaFilhoGouveia2003} apresentam os SSD como instrumentos capazes de integrar conhecimento técnico ao planejamento, em que a simulação responde a perguntas do tipo ``o que acontece se'' e a apresentação de resultados por gráficos, tabelas e indicadores facilita a comparação entre alternativas.

\section{Ferramentas de simulação aplicadas a sistemas de reservatórios}
\label{sec:trabalhos-relacionados}

Diversas ferramentas computacionais foram desenvolvidas para apoiar a análise da operação de sistemas de reservatórios. O MODSIM, desenvolvido na Colorado State University, é um sistema de suporte à decisão para gestão de bacias que representa o sistema como uma rede de fluxo e resolve, a cada intervalo, um problema de otimização que distribui a água segundo prioridades definidas pelo usuário \cite{labadie2004,labadie2007modsim}. No Brasil, o AcquaNet, desenvolvido pelo Laboratório de Sistemas de Suporte a Decisões da Escola Politécnica da Universidade de São Paulo com base no MODSIM, adota a mesma formulação de rede de fluxo e é empregado em estudos de alocação de água em sistemas complexos \cite{schardong2009acquanet}. O HEC-ResSim, do Corpo de Engenheiros do Exército dos Estados Unidos, simula a operação de reservatórios por regras, com intervalos de tempo que podem ser diários ou menores e representação detalhada de zonas operacionais e restrições \cite{klipsch2013ressim}. O WEAP integra oferta e demanda de água no planejamento de bacias, com alocação orientada por prioridades de demanda e preferências de fonte \cite{yates2005weap}, e o RiverWare oferece um ambiente genérico para construir modelos detalhados de sistemas de reservatórios, com simulação por regras e otimização \cite{zagona2001riverware}.

Essas ferramentas são genéricas e abrangentes: podem representar redes extensas, intervalos de tempo curtos e diferentes objetivos operacionais. Em contrapartida, cada aplicação exige que o usuário construa a rede, carregue as séries hidrológicas e as relações físicas dos reservatórios e traduza as regras de operação para a linguagem do programa, o que demanda tempo e conhecimento especializado. O Quadro \ref{qua:ferramentas} resume as características dessas ferramentas e as do sistema desenvolvido neste trabalho.

\begin{quadro}[!htbp]
    \captionsetup{width=15cm}
    \caption{Características de ferramentas de simulação de sistemas de reservatórios e do sistema desenvolvido}
    \label{qua:ferramentas}
    \centering
    \footnotesize
    \renewcommand{\arraystretch}{1.35}
    \begin{tabular}{|>{\raggedright\arraybackslash}p{2.1cm}|>{\raggedright\arraybackslash}p{3.4cm}|>{\raggedright\arraybackslash}p{2.0cm}|>{\raggedright\arraybackslash}p{2.4cm}|>{\raggedright\arraybackslash}p{3.6cm}|}
        \hline
        \textbf{Ferramenta} & \textbf{Abordagem} & \textbf{Intervalo} & \textbf{Acesso} & \textbf{Dados e regras} \\
        \hline
        MODSIM & rede de fluxo com otimização a cada intervalo & definido pelo usuário & programa instalado & rede, séries e prioridades montadas pelo usuário \\
        \hline
        AcquaNet & rede de fluxo com otimização, derivado do MODSIM & mensal, entre outros & programa instalado & rede, séries e prioridades montadas pelo usuário \\
        \hline
        HEC-ResSim & simulação por regras e zonas operacionais & diário ou menor & programa instalado & reservatórios, regras e séries configurados pelo usuário \\
        \hline
        WEAP & planejamento integrado de oferta e demanda & mensal, entre outros & programa instalado, mediante licença & cenários, demandas e prioridades definidos pelo usuário \\
        \hline
        RiverWare & simulação por regras e otimização & definido pelo usuário & programa instalado, mediante licença & modelo construído pelo usuário \\
        \hline
        Sistema desenvolvido & simulação do balanço hídrico mensal com regras por estado operacional & mensal & navegador web & base com 155 açudes do Ceará, hidrossistemas e curvas-guia dos planos de gestão proativa de secas já cadastrados \\
        \hline
    \end{tabular}
    \Fonte{Elaborada pelo autor com base em \citeonline{labadie2007modsim}, \citeonline{schardong2009acquanet}, \citeonline{klipsch2013ressim}, \citeonline{yates2005weap} e \citeonline{zagona2001riverware}.}
\end{quadro}

Uma alternativa mais simples a essas ferramentas é o cálculo do balanço hídrico em planilha eletrônica. A planilha é flexível e transparente, mas cada estudo exige preparar novamente séries, curvas e regras, cada cenário demanda ajustes manuais nas fórmulas, e a revisão do cálculo torna-se difícil à medida que aumentam o número de reservatórios, as regras de operação e os cenários analisados. Em relação a ela, o sistema desenvolvido mantém a transparência do cálculo mensal, tanto que foi verificado contra uma planilha independente, e acrescenta a base de dados já organizada, a aplicação automática das regras por estado operacional, a operação conjunta de reservatórios e a execução de novos cenários em poucos segundos. O sistema desenvolvido neste trabalho não pretende substituir as ferramentas descritas nesta seção, que são mais completas, por exemplo na otimização da alocação e na representação de intervalos curtos e redes extensas. Seu diferencial está na combinação de características voltadas ao contexto cearense: base de dados já organizada com os reservatórios monitorados no estado, representação direta das regras por estados operacionais adotadas nos planos de gestão proativa de secas, acesso pelo navegador, sem instalação, e disponibilização dos registros mensais completos de cada simulação, acompanhada de verificação documentada. Essa combinação reduz o esforço necessário para examinar cenários de operação, que é a finalidade de um sistema de suporte à decisão segundo \citeonline{sprague1980}.

\section{Verificação de modelos computacionais}
\label{sec:verificacao-modelos}

A literatura de modelagem distingue verificação e validação. A verificação consiste em demonstrar que o programa computacional e sua implementação estão corretos, isto é, que o código resolve as equações e executa as regras do modelo conceitual conforme especificado. A validação consiste em demonstrar que o modelo, dentro de seu domínio de aplicabilidade, representa o sistema real com exatidão suficiente para o uso pretendido \cite{refsgaardHenriksen2004,sargent2013}. De forma resumida, a verificação pergunta se as equações foram resolvidas corretamente; a validação pergunta se as equações corretas foram escolhidas \cite{oberkampfRoy2010}.

As duas atividades exigem evidências diferentes. A verificação pode ser realizada sem dados observados, por meio da comparação com cálculos independentes, do controle da conservação de massa, de testes em casos controlados com resposta conhecida e da inspeção de condições de contorno, como limites físicos e mudanças de estado. A validação exige confronto com observações do sistema real, como volumes medidos nos reservatórios, preferencialmente em períodos não utilizados para ajustar o modelo \cite{oberkampfRoy2010,refsgaardHenriksen2004}.

A análise de sensibilidade complementa a verificação ao examinar como as saídas respondem a variações nas entradas. Quando a resposta tem a direção e a ordem de grandeza esperadas, aumenta-se a confiança de que o modelo se comporta de forma coerente; quando não tem, a análise pode revelar erros de implementação ou hipóteses inadequadas. A variação de um parâmetro por vez, mantendo os demais fixos, é uma forma simples e transparente de análise local, embora não capture interações entre parâmetros \cite{pianosi2016}.
